% Guide (unnumbered). The whole story as one carved scene, read bottom to top: the configuration space C as a
% hatched slab with the hole the carving leaves and one configuration as a point of it; the feasible solid F lifted
% out and cut open, with the two generator paths and the version points on its core sheet and a packet on one
% version; above, the quotient patch V with the two loops, and over V the six sheets, one per version, on which the two loops
% lift: neither lift closes, and they end on different sheets (monodromy); the sheets carry no names on the cover.
% The four parts of the paper are headings
% beside their objects, with section ranges only.
% The solid's decoration is turned by 15 degrees so that the projection line from H to [X] passes no other point.
% Colours (the anchored palette): red = versions, teal tint = the core sheet, gold = the
% packet, white = cladding, sheets and the quotient patch, blue = nitrogen. Ink carries the drawing; colour only guides the eye. The two actions are ink, t solid and tu dashed,
% on the solid and on the sheets alike. Hatch = cut material. Arrowheads only on the two paths on the solid.
\documentclass[10pt,border=1mm]{standalone}
\usepackage[T1]{fontenc}
\usepackage{figurestyle}
\input{primitives.tex}
\input{molecules.tex}
\input{numbers.tex}
\begin{document}
\begin{tikzpicture}[plate,head/.style={font=\fontsize{9.5}{11}\selectfont\bfseries,inner sep=0pt},rng/.style={font=\fontsize{8}{10}\selectfont,inner sep=0pt}]
\path[use as bounding box] (0,1.1) rectangle (15,15.6);
\pgfmathsetmacro{\Ro}{2.3}\pgfmathsetmacro{\Ri}{1.62}\pgfmathsetmacro{\E}{.42}
% ------------------------------------------------ the slab C, cut away everywhere, with the hole the solid leaves
\path[cut] (.5,1.8)--(12.1,1.8)--(14.5,4.0)--(2.9,4.0)--cycle; \draw[heavy] (.5,1.8)--(12.1,1.8)--(14.5,4.0)--(2.9,4.0)--cycle;
\node[lab,anchor=north west] at (.6,1.68) {$\conf$};
\path[fill=white,even odd rule] (7.6,2.9) ellipse[x radius=\Ro,y radius={\E*\Ro}] (7.6,2.9) ellipse[x radius=\Ri,y radius={\E*\Ri}];
\draw[fine] (7.6,2.9) ellipse[x radius=\Ro,y radius={\E*\Ro}]; \draw[fine] (7.6,2.9) ellipse[x radius=\Ri,y radius={\E*\Ri}];
% ------------------------------------------------ one configuration, a point of C
\fill[PlateInk] (4.0,3.25) circle[radius=1.4pt];
\begin{scope}\molsolid\mollabelsfalse\pic[scale=.72] at (4.0,6.1) {mol3d-mla-ref};\end{scope}
\draw[leader] (4.0,5.28)--(4.0,3.31);
\node[head,anchor=west] at (.4,4.95) {I\quad Configurations}; \node[head,anchor=west] at (.4,4.55) {and States}; \node[rng,anchor=north west] at (.4,4.33) {\S\S\,1--3};
% ------------------------------------------------ the feasible solid F lifted out: core sheet, two paths, versions, a packet
\begin{scope}[shift={(7.6,6.2)}]\tikzset{tube/Ro=2.3,tube/Ri=1.62,tube/H=.55,tube/E=.42,tube/a1=-47,tube/a2=23}
  % the decoration is turned by phi = 15 degrees so that no version point shares a vertical with its antipode
  \pgfmathsetmacro{\ph}{15}
  \draw[hidden] (-\Ro,-3.3)--(-\Ro,-.55) (\Ro,-3.3)--(\Ro,-.55);
  \pic (ft) at (0,0) {thicktube};
  \foreach \a in {0,120,240} {\coordinate (vt\a) at ({1.96*sin(\a+\ph)},{.55-.42*1.96*cos(\a+\ph)});}
  \foreach \a in {60,180,300} {\coordinate (vb\a) at ({1.96*sin(\a+\ph)},{-.55-.42*1.96*cos(\a+\ph)});}
  \draw[tpath] plot[domain=\ph:\ph+115,samples=30,variable=\a] ({1.96*sin(\a)},{.55-.42*1.96*cos(\a)});
  \draw[upath] plot[domain=0:1,samples=30,variable=\s] ({1.96*sin(\ph-60*\s)},{.55-2*.55*\s-.42*1.96*cos(\ph-60*\s)});
  \path[fill=ColRetained!55] (vt0) circle[radius=.3]; \draw[fine] (vt0) circle[radius=.3];
  \foreach \p in {vt0,vt120,vt240,vb300} {\hatomat{\p}{.085}{ColDot}}
  \foreach \p in {vb60,vb180} {\draw[ghost,fill=ColDot!45] (\p) circle[radius=.085];}
  \node[lab,anchor=east] at ({-\Ro-.18},-1.1) {$\mathcal F$};
  \coordinate (Hpt) at (vt0);
\end{scope}
\node[head,anchor=west] at (10.7,6.1) {II\quad Fragments and}; \node[head,anchor=west] at (10.7,5.7) {Feasible Regimes}; \node[rng,anchor=north west] at (10.7,5.48) {\S\S\,4--6};
\node[head,anchor=west] at (.5,8.9) {III\quad Symmetry and}; \node[head,anchor=west] at (.5,8.5) {Localization}; \node[rng,anchor=north west] at (.5,8.28) {\S\S\,7--9};
% ------------------------------------------------ the quotient patch V with the two loops; over it the six sheets with the two lifts
\draw[hidden] (Hpt)--(8.1,6.08); \draw[hidden] (8.1,7.72)--(8.11,9.52);
\node[sub,anchor=west,inner sep=1pt] at (8.25,8.25) {$\mathcal F\to\mathcal F/G_{12}$};
\begin{scope}[shift={(7.76,9.3)}]
  \newcommand\sheet[3]{\path[fill=#3] (-1.8,#1)--(1.8,#1)--(2.5,#1+.5)--(-1.1,#1+.5)--cycle; \draw[#2] (-1.8,#1)--(1.8,#1)--(2.5,#1+.5)--(-1.1,#1+.5)--cycle;}
  \sheet{0}{heavy}{white}
  \coordinate (bx) at (.35,.25); \fill[PlateInk] (bx) circle[radius=1.2pt];
  \draw[tlift] (bx) .. controls (2.2,-.22) and (2.2,.72) .. (bx);
  \draw[ulift] (bx) .. controls (-1.5,.72) and (-1.5,-.22) .. (bx);
  \foreach \k/\la in {0/H,1/uH,2/tH,3/tuH,4/{t^2H},5/{t^2uH}} {\pgfmathsetmacro{\yy}{1.3+.78*\k}
    \sheet{\yy}{fine}{white}
    \coordinate (p\k) at (.35,{\yy+.25});
    }
  % the same two loops lifted from the sheet H: t climbs to tH, tu to tuH; neither closes, and they end on different sheets
  \draw[tlift] (p0) .. controls (2.3,{1.3+.25+.3}) and (2.3,{1.3+.78*2+.25-.3}) .. (p2);
  \draw[ulift] (p0) .. controls (-1.7,{1.3+.25+.4}) and (-1.7,{1.3+.78*3+.25-.4}) .. (p3);
  \foreach \k in {0,...,5} {\hatomat{p\k}{.075}{ColDot}}
\end{scope}
\node[head,anchor=west] at (10.7,13.3) {IV\quad Representations and}; \node[head,anchor=west] at (10.7,12.9) {Global Structure}; \node[rng,anchor=north west] at (10.7,12.68) {\S\S\,10--12};
% ------------------------------------------------ the symmetry groups: the bond interval [H, B], grey, covers in the line grammar
\def\gt{t}\def\gu{u}
\begin{scope}[shift={(3.2,10.0)},x=-0.5cm,y=1.4cm]
  \foreach \i/\x/\y/\v/\name in \hasseNodes {\coordinate (h\i) at (\x,\y);}
  \foreach \a/\b/\g in \hasseCoverLabels {\edef\gg{\g}\ifx\gg\gt \draw[tlift,draw=PlateInk!55] (h\a)--(h\b);\else\ifx\gg\gu \draw[ulift,draw=PlateInk!55] (h\a)--(h\b);\else \draw[blift,draw=PlateInk!55] (h\a)--(h\b);\fi\fi}
  \draw[tlift] (h0)--(h4); \draw[ulift] (h4)--(h6); \draw[blift] (h6)--(h9);
  \foreach \i/\x/\y/\v/\name in \hasseNodes {\draw[line width=.5pt,draw=PlateInk!55,fill=white] (h\i) circle[radius=2pt];}
  \foreach \i in {0,4,6,9} {\draw[line width=.5pt,draw=PlateInk,fill=white] (h\i) circle[radius=2pt];}
  \node[sub,anchor=north,text=PlateInk] at ($(h0)+(0,-.12)$) {$H$}; \node[sub,anchor=south,text=PlateInk] at ($(h9)+(0,.12)$) {$B$};
\end{scope}
\end{tikzpicture}
\end{document}
